\section{Introduction} 
\label{sec:introduction}

A general, low-cost method for accurately simulating human behavior with AI agents would have wide application in the social sciences \citep{Charness2025NextGen,jackson2025ai,promising2025LLM}.
Recognizing this potential, a growing literature explores whether large language models (LLMs) can simulate human responses in various settings.\footnote{\citep{argyle2022out,aher2022using,binz2023cognitive,brand2023using,park2023generative,
jackson2024turing,1000agents2024park,chang2024networks,Manning2024Automated,
capra2024llms,hansen2024simulating,kim2024homoeconomicus,
wang2025market,Duch2025Possum,zhu2025evidence,Broska2025Mixed,
bybee2025ghost,fish2025collusion}}
Across dozens of experiments, samples of these agents respond with remarkable similarity to humans---even when simulating novel studies that did not appear in the underlying LLM's training corpus \citep{hewitt2024predictingLLM,binz2025centaur,ValidityLLM2024,Tranchero2024LLMtheorizing,suh2025finetuningscaled}.
Yet within this literature, others find settings where the very same models are poor human proxies.\footnote{\citep{opinions2023Santukar,WhichHumans:2023,Cheng2023caricature,Gui2023Causal,gao2024caution}}
This inconsistency poses a challenge for AI simulations as robust and credible predictive models---particularly in settings where no prior human data exists.
The core challenge is not simply achieving a close match between AI and human responses in one setting, but building agents that will generalize reliably.

A natural starting point is improving the instructions given to agents.
These instructions, or ``\persona{s},'' are written descriptions given to the LLM specifying who it is, what it believes, or how it should behave and reason \citep{horton2023large}.
Such second-person instructions (e.g., \emph{``You respond as type-X person''}) can profoundly affect output distributions because advanced LLMs have been explicitly fine-tuned to follow instructions \citep{bai2022helpfulharmless,ouyang2022training}.
With an appropriate prompt (which can be massive),\footnote{To date, Google's Gemini 1.5 can accommodate 10 million tokens \citep{Gemini152024}. 
This is roughly equivalent to 15{,}000 pages, or about 20 copies of the Handbook of Experimental Economics \citep{Kagel1995Handbook}.} highly capable models can perform complex reasoning and mathematical tasks at levels sometimes better even than those of highly skilled humans.

Despite the existence of these powerful, steerable models, constructing agents whose behavior is similar to that of real humans in a wide variety of settings is nontrivial---even when human data are available to guide the search.
The set of possible \persona{s} is vast, ranging from simple combinations of social or demographic traits to complex programmatic instructions related to how humans make decisions \citep{zhu2025ExplPredicting,Jackson2025Mixture}.
As in other machine learning applications, the challenge is not only to avoid underfitting but also to guard against overfitting.
By iterating through enough \persona{s}, one can almost always find some arbitrary \persona{} that shifts the LLM's responses to closely match a given human distribution.
For example, an LLM instructed \emph{``you randomly offer between \$6 and \$9''} may perfectly reproduce a distribution of human responses in a \$20 dictator game, but such a \persona{} would be nonsensical for a \$5 dictator game.
In contrast, a \persona{} grounded in the underlying behavioral drivers---e.g., \emph{``you are self-interested but fair''}---can perform well in-sample and plausibly extend to a range of allocation games.
Standard data-driven approaches, such as a train-test split within a single dataset, cannot reliably distinguish between these two cases; the latter appears better only when tested in truly new settings.
If the goal is to predict behavior in settings with no prior human data, how should researchers construct \persona{s}?

In this paper, we build agents whose behavior in simulations usefully generalizes to what we see from humans across entire domains.
Our approach mirrors what researchers generally try to do in social science, but in reverse. 
Rather than testing a theory with empirical data, theory is embedded in agents (via natural language instructions), which are then used to generate candidate data.
The theory and agent composition is then optimized to reduce error with respect to real-world data from a domain where predictions are desired.
Human data from distinct but conceptually similar settings serve as held-out test sets or are incorporated into training to improve generalization.
We show that ``general'' agents constructed and validated in this way can improve the predictive power of \aisub{s} in novel settings.
Both steps are essential: without theoretical grounding, optimized \persona{s} may fail to meaningfully improve even in-sample predictions, and without cross-setting validation, they are prone to overfit.

The approach uses two kinds of data: (i) \emph{training data}---existing human data used to optimize AI simulations, and (ii) \emph{validation} or \emph{test data}---existing human data related to but distinct from the training data, used to test whether the optimized agents generalize.
The ultimate goal is to produce agents that can more accurately predict human behavior in novel \emph{target settings} where no prior human data exist, but that lie in the same broad domain as the training and validation settings.

The first step of the approach is to limit the ``space'' of \persona{s} to a subset motivated by some economic theory or causal mechanism relevant to the novel setting of interest.
This theory-grounding is analogous to constraining the functional form of the hypothesis class in machine learning.
Continuing with the dictator game example, where social preferences likely determine behavior, one might choose the candidate set of \persona{s} characterized by known drivers of the relevant preferences (e.g., \emph{``You are \{level\}}'' for all \emph{levels} $\in$ \{\emph{self-interested but fair, altruistic, selfish}\}).

The second step is to optimize over this set to best match the human training data.\footnote{There is an active literature on optimizing prompts \citep{khattab2024dspy}.}
All candidate \persona{s} are put in simulations of the settings that produced the training data (e.g., a \$20 dictator game).
Optimization effectively filters these candidate \persona{s} down to a final subset whose simulated responses are closest to the human training distribution.
We employ two optimization methods: (i) a selection method that identifies the optimal mixture of \persona{s} from the candidate set \citep{leng2024calibrate,Jackson2025Mixture,bui2025mixture}, and (ii) a construction method that optimizes numerical parameters embedded directly in the prompts.
If the final optimized set achieves a good in-sample fit, we have reason to believe they will generalize because they are also grounded in a relevant theory.
Poor in-sample fit suggests a mismatch between the theory and the training setting or an inadequate operationalization of that theory.
In this case, we revise the candidate \persona{s} and re-optimize.

Given strong in-sample fit, we assess whether the final set of \persona{s} generalize using a train-test split approach inspired by the principles of invariant risk minimization \citep{Peters2016invariant,Heinze2018invariant,arjovsky2020invariant}.
The final set of \persona{s} is placed in simulations of the settings that produced the human validation data, where we also expect the underlying theory to hold (e.g., optimize on a \$20 dictator game, but test on a \$5 dictator game).
To be clear, this means the validation set necessarily comes from a distinct data-generating process from the one that produced the training data. 
We then compare the predictions from these simulations to the relevant human distributions.
By construction, sets of \persona{s} with strong test performance---those that accurately predict the validation data---are then those that capture generalizable relationships predictive of human behavior across contexts.
Those that fail validation likely do not.
Consequently, if the novel target setting is governed by the same theory or causal mechanism used to construct the optimized \persona{s} (e.g., the target is a \$50 dictator game), we may gain confidence that they will better predict human responses in that setting.

We illustrate this approach and provide evidence of its efficacy with training and validation data drawn from experiments in the behavioral economics literature.
All simulations use \textsc{Gpt-4o} with the temperature set to 1, though the approach is agnostic to the choice of model and hyperparameters.\footnote{\textsc{Gpt-4o}, among the most widely used LLMs, and the temperature setting are defaults for the software used to run our simulations \citep{Horton2024EDSL}.}
We first apply the selection method to \citet{1120Arad2012}'s 11-20 money request game, where participants request an amount and receive a bonus if they choose exactly one less than their opponent. 
We use only the original dataset from the paper, which contained fewer than 200 observations.
This setting is appealing to study because optimal play depends not only on the focal agent's capabilities, but also on their beliefs about how others will reason.
Endowing \aisub{s} with distinct \persona{s} corresponding to varying degrees of strategic reasoning (the theoretical focus of \citeauthor{1120Arad2012}) produces a mixture that closely matches the original human data.

When we validate these optimized samples on distinct variants of the 11-20 money request game, they are better predictors of initial human play than baseline \aisub{s} with no additional instructions.
By contrast, scientifically meaningless or ``atheoretical'' \aisub{s} derived from historical figures, pseudo-scientific Myers-Briggs personality types, and those instructed to select particular numbers can sometimes match human distributions in one variant of the game but fail to generalize across others.

We next test the predictive power of the optimized agents in target settings where no prior human data exist.
To do so, we construct four new games and collect responses from samples of preregistered crowdsourced participants \citep{horton2011online} on Prolific. 
These games are derived from the original 11-20 game (and its variants), but adapted to other numeric ranges (1-10 and 1-7).
The optimized sample of theory-grounded \persona{s} produces responses that predict the new human data far better than the off-the-shelf baseline.
Prediction error is decreased by \MinPercImproveArad\%-\MaxPercImproveArad\% across the games.
Furthermore, these simulations predict the results of the new experiments in some games \emph{better} than the most plausibly relevant human data from \citeauthor{1120Arad2012}; in one case, the KL divergence is halved.
By contrast, the alignment of the atheoretical \persona{s}---which failed validation---with the new human data is often similar to or worse than the baseline AI.

What statistical guarantees does this approach afford?
Without a correctly specified causal model, no statistical procedure can guarantee performance in arbitrary new environments \citep{Pearl2009Causality}.  
Formal guarantees with existing data, like those required for prediction-powered inference \citep{Angelopoulos2023Prediction} and other related methods \citep{Egami2023Using,Hardy2025Population}, would require a strictly firewalled validation set: data never used in the construction of the underlying LLM \citep{ludwig2025llmapplied,sarkar2024lookahead,ML2017Sendhil,spiess2025causaltext}.
This is a tall order impossible to meet in practice, even with existing public weight LLMs, never mind private models. 
However, what we can guarantee is prediction performance over a \emph{pre-committed} family of settings. 

The setup is very similar to the theoretical framework of \cite{andrews2025transfer}, who study how well predictions from different economic and machine learning models transfer across economic domains.
The first step is to establish a clearly defined space of experimental settings---for example, variants of public goods games, dictator games, or other allocation games differentiated by instructions, parameters, or other structural features.
This idea is related to that in the literature on program evaluation and external validity, where population-level treatment effects can be estimated by randomizing over a defined set of possible samples \citep{Hotz2005PredictingEfficacy,Allcott2015Site,Dehejia2019External}.
From this space, a random subset of settings is drawn and randomly assigned to human subjects, who provide responses.\footnote{This is also similar to the integrative experiment designs of \cite{20Q2024Almaatouq}.}
By comparing these human responses with LLM-generated predictions, we can make externally valid estimates across the entire space.

A key distinction from the setups laid out in \citeauthor{Hotz2005PredictingEfficacy}, \citeauthor{Allcott2015Site}, and \citeauthor{Dehejia2019External} is that appropriate coverage does not require that the family of settings share a data-generating process or that the underlying samples of human subjects are from the same population.
Indeed, the variance of estimates naturally reflects how closely the held-out settings resemble those sampled.  
If all scenarios are variants of a single setting---e.g., a dictator game with different monetary amounts---performance is tightly bounded; if they span disparate domains---e.g., many types of allocation games---estimates are likely less precise.
Note that even if a setting inadvertently appears in the model's pre-training data, the random sampling procedure still accounts for its contribution---such cases may merely reduce prediction error.

We explore this approach to inference by constructing a population of \TotalGames{} novel strategic games.
These were inspired by the original 11-20 game but were modified along several dimensions, such as support of possible choices, the size and nature of the bonus, the manner in which money is allocated, etc. 
The population of games is the full factorial combination of these dimensions.
The resulting games differed substantively from each other; in some cases, they were cooperative, others zero-sum, many had a dominant strategy equilibrium, while some had no dominant strategies at all.
The vast majority are surely not in any LLM's training corpus to date.
From the population, we sample \NGamesS{} unique games and have \NParticipantS{} human subjects each play a game in a final preregistered experiment.
We take the theoretically-grounded level-$k$ agents---agents constructed months before the novel set of \TotalGames{} games existed---and evaluate their ability to predict the human responses.

The theory-grounded agents are significantly better predictors of initial human play than the baseline AI off-the-shelf.
In the average game, they assign \PctBaselineBarLambda\% more probability to the actions actually taken by human subjects.
The optimized agents are similarly effective compared to two well-established theoretical benchmarks: (i) a symmetric Nash equilibrium calculated for each of the 1,500 games (\PctNashBarLambda\% more probability), and (ii) the game-specific predictions from the cognitive hierarchy model of \cite{Camerer2004Cognitive} (\PctCHBarLambda\% more probability).
Because the games were randomly sampled, the corresponding confidence intervals over the relative predictive power are externally valid for the much larger population.

The goal of AI simulations in this paper is to harness two extensive sources of information to create better predictive agents: (i) well-established theoretical models from the social sciences, and (ii) the immense knowledge about human behavior that LLMs have implicitly learned during training \citep{lindsey2025biology,ameisen2025circuit}.
In a sense, \aisub{s} are a general vessel through which theory can be flexibly applied to any setting.
Our approach aims to generate such agents by conducting a structured trial run across various settings, building evidence that theoretically-grounded \persona{s} generalize effectively to similar yet distinct environments.
These principles can be applied to any domain where relevant human data exists---even those with multiple interacting agents \citep{Manning2024Automated,qian2025strategicbargaining}.
One can imagine iteratively identifying novel scenarios where predictions are needed, finding the closest relevant existing data, and continuously optimizing and calibrating agents as needed.\footnote{LLMs could automatically construct candidate \persona{s} as input to such feedback loops \citep{Jackson2025Mixture}.}
Our main contribution is to systematize these ideas for constructing agents and then provide empirical evidence that they can substantially improve the predictive power of AI simulations in new settings.

The remainder of this paper is organized as follows.
Section~\ref{sec:approach} begins with a concrete example of identifying generalizable relationships and constructing \persona{s} that better predict human subjects in new settings.
We then describe the approach more generally and specify the optimization methods.
Section~\ref{sec:predict_new_AR} illustrates the approach and demonstrates its efficacy empirically with several experiments.
Section~\ref{sec:ext-valid} demonstrates how agents can be used in a pre-committed setting to provide externally valid statistical estimates of predictive accuracy at scale.
The paper concludes in Section~\ref{sec:conclusion}.
